%====================================================
% 02 Objectives
%====================================================

\section{Objectives}

The primary objective of this assignment is to develop and evaluate a Convolutional Neural Network (CNN)-based handwritten digit recognition system using the MNIST handwritten digit dataset. The assignment also aims to investigate how different CNN architectural configurations influence classification performance and computational efficiency.

The specific objectives are as follows:

\begin{itemize}

\item To understand the characteristics and structure of the MNIST handwritten digit dataset.

\item To preprocess the dataset by normalizing pixel values and reshaping the images into the format required for CNN training.

\item To design and implement multiple CNN architectures using TensorFlow and Keras.

\item To investigate the effect of different convolution kernel sizes on handwritten digit classification performance.

\item To evaluate the influence of increasing the number of convolution layers within the CNN architecture.

\item To compare the performance of different optimization algorithms, namely Adam and Stochastic Gradient Descent (SGD).

\item To investigate the effect of different activation functions, including ReLU and Tanh.

\item To train and validate each CNN model using the MNIST training dataset.

\item To evaluate the performance of each model using testing accuracy, testing loss, confusion matrix analysis, and misclassification analysis.

\item To compare all developed CNN models based on classification accuracy, computational complexity, training time, and prediction performance.

\item To validate the trained CNN using custom handwritten digit images created independently of the MNIST dataset.

\item To identify the most effective CNN architecture for handwritten digit recognition based on the experimental results.

\item To gain practical experience in implementing, training, evaluating, and comparing deep learning models for image classification applications.

\end{itemize}